%-------------------------
% Resume in Latex
% Author : Pranav Damera
%------------------------

\documentclass[letterpaper,11pt]{article}

\usepackage{latexsym}
\usepackage[empty]{fullpage}
\usepackage{titlesec}
\usepackage{marvosym}
\usepackage[usenames,dvipsnames]{color}
\usepackage{verbatim}
\usepackage{enumitem}
\usepackage[hidelinks]{hyperref}
\usepackage{fancyhdr}
\usepackage[english]{babel}
\usepackage[utf8]{inputenc}
\usepackage{newunicodechar}
\newunicodechar{₂}{$_2$}
\usepackage{tabularx}
\usepackage{fontawesome5}
\usepackage{multicol}
\setlength{\multicolsep}{-3.0pt}
\setlength{\columnsep}{-1pt}
\input{glyphtounicode}

\pagestyle{fancy}
\fancyhf{}
\fancyfoot{}
\renewcommand{\headrulewidth}{0pt}
\renewcommand{\footrulewidth}{0pt}

\addtolength{\oddsidemargin}{-0.6in}
\addtolength{\evensidemargin}{-0.5in}
\addtolength{\textwidth}{1.19in}
\addtolength{\topmargin}{-.7in}
\addtolength{\textheight}{1.4in}

\urlstyle{same}
\raggedbottom
\raggedright
\setlength{\tabcolsep}{0in}

\titleformat{\section}{
  \vspace{-4pt}\scshape\raggedright\large\bfseries
}{}{0em}{}[\color{black}\titlerule \vspace{-5pt}]

\pdfgentounicode=1

\newcommand{\resumeItem}[1]{
  \item\small{
    {#1 \vspace{-2pt}}
  }
}

\newcommand{\resumeSubheading}[4]{
  \vspace{-2pt}\item
    \begin{tabular*}{1.0\textwidth}[t]{l@{\extracolsep{\fill}}r}
      \textbf{#1} & \textbf{\small #2} \\
      \textit{\small#3} & \textit{\small #4} \\
    \end{tabular*}\vspace{-7pt}
}

\newcommand{\resumeProjectHeading}[2]{
    \item
    \begin{tabular*}{1.001\textwidth}{l@{\extracolsep{\fill}}r}
      \small#1 & \textbf{\small #2}\\
    \end{tabular*}\vspace{-7pt}
}

\renewcommand\labelitemi{$\vcenter{\hbox{\tiny$\bullet$}}$}
\renewcommand\labelitemii{$\vcenter{\hbox{\tiny$\bullet$}}$}

\newcommand{\resumeSubHeadingListStart}{\begin{itemize}[leftmargin=0.0in, label={}]}
\newcommand{\resumeSubHeadingListEnd}{\end{itemize}}
\newcommand{\resumeItemListStart}{\begin{itemize}}
\newcommand{\resumeItemListEnd}{\end{itemize}\vspace{-5pt}}

%-------------------------------------------
\begin{document}

%----------HEADING----------
\begin{center}
    {\Huge \scshape Pranav Damera} \\ \vspace{1pt}
    Aldie, VA 20105 \\ \vspace{1pt}
    \small \raisebox{-0.1\height}\faPhone\ 571-278-6512 ~
    \href{mailto:pdamera05@gmail.com}{\raisebox{-0.2\height}\faEnvelope\ \underline{pdamera05@gmail.com}} ~
    \href{https://github.com/pranavdamera}{\raisebox{-0.2\height}\faGithub\ \underline{github.com/pranavdamera}} ~
    \href{https://linkedin.com/in/psdamera}{\raisebox{-0.2\height}\faLinkedin\ \underline{linkedin.com/in/psdamera}}
    \vspace{-8pt}
\end{center}

%-----------EDUCATION-----------
\section{Education}
  \resumeSubHeadingListStart
    \resumeSubheading
      {George Mason University}{Aug 2023 -- May 2028}
      {Bachelor of Science in Computer Engineering}{Fairfax, Virginia}
  \resumeSubHeadingListEnd

%-----------EXPERIENCE-----------
\section{Experience}
  \resumeSubHeadingListStart

    \resumeSubheading
      {iTech AG}{June 2026 -- Aug 2026}
      {AI Foundry Intern}{Arlington, Virginia}
      \resumeItemListStart
        \resumeItem{Built agentic AI workflows on \textbf{Amazon Bedrock AgentCore}, orchestrating tool-calling LLM agents to automate [DESCRIBE TASK, e.g. candidate/client intake triage] end-to-end, cutting manual processing time by [X]\% across a government/commercial IT consulting pipeline.}
        \resumeItem{Made hands-on infrastructure decisions provisioning and schema-designing \textbf{Amazon Aurora PostgreSQL} as the system of record for project, candidate, and client engagement data -- owning indexing and connection-pooling strategy to support [X] concurrent users at sub-[X]ms query latency.}
        \resumeItem{Architected a \textbf{SharePoint}-to-Aurora integration (Microsoft Graph API) that structured unstructured lists/documents into relational schemas feeding a REST API and dashboard, eliminating [X] hours/week of manual data reconciliation and improving resource-allocation visibility across active delivery pipelines.}
      \resumeItemListEnd


    \resumeSubheading
      {GoldenFlow Labs}{Sep 2025 -- Present}
      {Founder \& ML Systems Engineer}{Arlington, Virginia}
      \resumeItemListStart
        \resumeItem{Designed a FastAPI NIR spectral classification service for physical product authentication, processing structured 900--1700nm spectral arrays through preprocessing, CNN-based inference, and structured result storage.}
        \resumeItem{Architected a 1D-CNN with attention for spectral feature extraction and risk classification, applying convolutional feature learning concepts across structured sensor arrays; exported lightweight inference artifacts through ONNX for reproducible deployment workflows via Git-tracked CI/CD pipelines.}
        \resumeItem{Implemented tamper-evident audit architecture linking inference outputs to cryptographic hashes, enabling reproducible traceability across physical product verification events with zero post-hoc modification.}
      \resumeItemListEnd

    \resumeSubheading
      {Avhita Health Systems}{Jun 2024 -- Dec 2025}
      {Software Engineering Intern}{Remote}
      \resumeItemListStart
        \resumeItem{Designed and deployed a real-time physiological data ingestion pipeline using FastAPI and Google Cloud Run, enabling continuous ECG streaming in 10--20 second windows with sub-2-second end-to-end latency across concurrent sessions.}
        \resumeItem{Implemented signal processing and validation workflows for ECG sensor streams at 250--500 Hz using Pan-Tompkins R-peak detection, reducing malformed data events by 40\% and improving arrhythmia detection reliability.}
        \resumeItem{Architected a multi-layer patient data platform integrating continuous sensor monitoring, an EHR data model, and longitudinal cardiac event tracking across patient cohorts with structured BLE device-to-cloud pipelines.}
      \resumeItemListEnd



  \resumeSubHeadingListEnd
\vspace{-14pt}

%-----------PROJECTS-----------
\section{Projects}
    \vspace{-5pt}
    \resumeSubHeadingListStart

      \resumeProjectHeading
          {\textbf{GridGuard -- Spatial Intelligence \& Anomaly Detection for Distributed Solar} $|$ \emph{Python, FastAPI, scikit-learn, XGBoost, pvlib, Next.js, MapLibre GL} $|$ \href{https://github.com/pranavdamera/gridguard}{\underline{GitHub}}}{}
          \resumeItemListStart
            \resumeItem{Built an open-source spatial intelligence platform that forecasts expected photovoltaic output and wraps it in a calibrated uncertainty band via \textbf{split conformal prediction} ($\alpha$=0.05, distribution-free) -- validated on real measured telemetry from 3 NIST/NREL PVDAQ arrays (586 kW combined), achieving R\textsuperscript{2} up to 0.993 and 96.2--99.1\% empirical coverage.}
            \resumeItem{Designed a physics-informed ML hybrid (pvlib PVWatts chain + gradient-boosted residual model) that outperformed pure ML on 2 of 3 real arrays, and a haversine-based geospatial neighborhood algorithm across a 7-site fleet to separate real equipment faults from shared weather effects via cross-site residual correlation.}
            \resumeItem{Engineered a reproducible offline-train / online-serve architecture (FastAPI backend, Next.js + MapLibre GL fleet-map frontend) where every served metric traces to a build manifest recording git commit, dataset window, and hyperparameters; enforced by 224 automated tests and CI across Python 3.11/3.12, reaching 100\% event-level fault recall (5/5) on real telemetry.}
          \resumeItemListEnd
          \vspace{-13pt}

      \resumeProjectHeading
          {\textbf{Dhara Energy Map -- Geospatial Intelligence Platform} $|$ \emph{Python, PostGIS, FastAPI, GeoJSON} $|$ \href{https://github.com/pranavdamera/dhara-energy-map}{\underline{GitHub}}}{}
          \resumeItemListStart
            \resumeItem{Built a geospatial intelligence platform for screening distributed energy resource sites using candidate parcel grids, solar potential, infrastructure proximity, land attributes, and observation snapshots.}
            \resumeItem{Designed PostGIS schemas and spatial indexes for candidate sites, infrastructure assets, observation snapshots, and transparent score breakdowns across ranked locations.}
            \resumeItem{Developed FastAPI services exposing GeoJSON site geometries, ranked site lists, parcel dossiers, and observation snapshots, with a data model designed for future raster and satellite monitoring overlays.}
          \resumeItemListEnd

      % -----------------------------------------------------------------------
      % THIRD PROJECT SLOT -- fill once built. Strongest candidate: a computer-
      % vision / satellite-imagery project, since GridGuard + Dhara already cover
      % telemetry anomaly detection and geospatial siting, but neither touches
      % imagery/CV. This is the differentiated gap for geospatial-intel /
      % defense recruiters (BlackSky, Anduril, Maxar, Planet Labs-style work).
      % TEMPLATE:
      %
      % \resumeProjectHeading
      %     {\textbf{Satellite Change Detection -- Sentinel-2 Land Cover Classification}
      %      $|$ \emph{Python, PyTorch, rasterio, scikit-learn}
      %      $|$ \href{https://github.com/pranavdamera/sentinel-change-detect}{\underline{GitHub}}}{}
      %     \resumeItemListStart
      %       \resumeItem{Fine-tuned a pretrained ResNet-18 CNN in PyTorch on multispectral Sentinel-2 satellite imagery
      %                   to classify land cover and detect temporal change across agricultural/infrastructure regions,
      %                   achieving [X]\% accuracy on a held-out validation set.}
      %       \resumeItem{Built a multi-image temporal processing pipeline ingesting paired before/after Sentinel-2 scenes,
      %                   computing NDVI difference rasters, and feeding stacked multi-band arrays into a CNN for
      %                   change-detection inference -- directly analogous to multi-image temporal analysis used in
      %                   commercial satellite-intelligence products.}
      %       \resumeItem{Deployed inference via a FastAPI endpoint with Dockerized runtime; outputs GeoJSON change masks
      %                   overlaid on an interactive map dashboard for visual validation of model predictions.}
      %     \resumeItemListEnd
      %     \vspace{-13pt}
      % -----------------------------------------------------------------------

    \resumeSubHeadingListEnd
\vspace{1pt}

%-----------TECHNICAL SKILLS-----------
\section{Technical Skills}
 \begin{itemize}[leftmargin=0.15in, label={}]
    \small{\item{
    \textbf{Languages}{: Python, Java, JavaScript, C/C++, SQL} \\
    \textbf{ML/AI}{: PyTorch, TensorFlow, scikit-learn, XGBoost, ONNX, CNNs, LSTM, multi-head attention, conformal prediction \& uncertainty quantification, physics-informed ML, Pandas, NumPy} \\
    \textbf{Geospatial \& CV}{: PostGIS, GeoJSON, rasterio, spatial indexing, geospatial neighborhood analysis (haversine), earth observation data, image-pair change detection} \\
    \textbf{Full-Stack}{: FastAPI, Flask, REST APIs, Next.js, React, MapLibre GL, Node.js, PostgreSQL, DynamoDB} \\
    \textbf{Cloud, Agentic AI \& DevOps}{: AWS (Bedrock, Bedrock AgentCore, Aurora PostgreSQL), GCP (Cloud Run), SharePoint/Graph API integration, LLM agent orchestration, Docker, Linux, CI/CD, Jenkins}
    }}
 \end{itemize}
 \vspace{-16pt}

\end{document}
